\chapter{설계 요구사항과 판정 기준}

\section{요구사항}

요구사항은 2026년 9월에 확정한 요구에 개선 구조와 검증에서 도출한 하위 요구를 더한 것이며, 값은 최종 설계(A25R)의 것이다(표 \ref{t:req}). 400±8 V 정밀 조정은 외부 부하의 요구가 아니라 고정이득 Stage C와 B\textsubscript{420} 부품 호환에서 나오는 파생 요구였다\cite{b9}, \cite{b10}. R2에서는 외부 요구인 버스 창 조건(부하 하한 이상, 부품 등급 상한 이하)을 직접 기준으로 쓴다. ISO 21498-1의 하위 등급 상한(420·470 V)은 과도와 리플을 제외한 정상 운전 전압으로 부품을 분류하는 값이다\cite{b9}. 이 과제는 전압 노드별로 상한을 나눠 적용한다. 400 V급 부품이 연결된 부하 버스에는 등급 상한(420·470 V)을 고장·충전 중 정지 같은 사건 중에도 적용하고(보수적인 설계 계약), 충전기 접속 노드 V\textsubscript{A}의 사건 한도는 사양 팩 클램프 시험 470 V, 고전압 W460 전단 시험 530 V다(W460 정상 최대 반전압 504 V, 결합 모서리 사건 V\textsubscript{A} 최대 516.37 V, E44). V\textsubscript{A} 결과를 창 조정단·보호 상태기계를 결합한 부하 버스 실증으로 넓히지 않으며, 부하 분리 이후의 보호 상태는 따로 판정한다. 이 상한들은 규격 요구나 부품의 과도 내량을 보인 것이 아니다(저전압 동작·부하 덤프·리플 내량은 따로 확인).

{\footnotesize
\begin{longtable}{L{\dimexpr 0.20\linewidth-2\tabcolsep\relax}L{\dimexpr 0.62\linewidth-2\tabcolsep\relax}L{\dimexpr 0.18\linewidth-2\tabcolsep\relax}}
\caption{설계 요구사항과 판정 기준}\label{t:req}\\
\toprule
\tabhead \textbf{항목} & \textbf{값} & \textbf{근거·상태} \\
\midrule
\endfirsthead
\toprule
\tabhead \textbf{항목} & \textbf{값} & \textbf{근거·상태} \\
\midrule
\endhead
\bottomrule
\endlastfoot
배터리 운전범위 & 사양 팩 640–920 V. 팩 등급 검토는 셀 180–240s(3.0–4.2 V) & E4, E18 \\
400 V 버스 창 & 부하 하한(260 V) 이상, 부품 등급 상한(B\textsubscript{420} 420 V, B\textsubscript{470} 470 V) 이하. 상한은 정상 전압 기준 하위 등급 값이고, 400 V급 부품이 연결된 부하 버스가 사건 중에도 넘지 않는 것은 이 과제의 설계 계약이다(V\textsubscript{A} 사건 한도는 V\textsubscript{A} 클램프 행). 창 조정 모듈은 V\textsubscript{W}(410·460 V) 기준 & 등급 값: 규격\cite{b9} / 사건 중 상한: 설계 계약 \\
창 조정 전압 품질 & Tier-1: V\textsubscript{W} ±8 V·리플 ≤ 0.1 V(리플은 PWM 주기평균 버스 전압의 마지막 1 ms 폭) / Tier-2: 상한 이하(주기평균 기준) & E13·E19 \\
연속전력·예산 & 25 kW(62.5 A), SALS 예산 23.75 kW & 부하 등급·계약 \\
전류 판정 기준 & 평균모델 Tier-1 90 A·Tier-2 93.75 A / 스위칭·회로 모델(리플 포함 순시) 98·100 A & 물리정격 100 A \\
제어 지연 & 모델 내 5 µs 영차유지 + 5 µs 전달지연, 모델 밖 예산 +3.0 µs & E6·E7 \\
창 조정 DAB 정격 & 정적 최대 링크전력 × 1.8(R2-W460 4.0 kW, R2-W410 5.0 kW) & 도출(E19) \\
센서·ADC(Stage B·창 조정 모듈) & 16비트, 범위 밖은 포화. 버스 0–600 V 단일 채널(잡음 표준편차 10 mV, 이득 오차 0.25 \%, 오프셋 1.2 V), 직렬 ±100 V·링크 0–150 V(5 mV), 전류 ±150 A(20 mA). 범위 600 V ≥ 1.1 × 사건 최대 V\textsubscript{A}·버스(516.3 V). 처음 구조 ③의 버스 350–450 V 차동 창(E12)은 400 V 정밀 조정 시험의 조건이다 & 도출(E12, E43) \\
게이트 시간 분해능 & H-브리지 PWM·DAB 위상 에지 ≤ 0.625 ns(고분해능 PWM) & 도출(E12, TB04B) \\
충전 시작 & 충전 전 Stage A 사전 기동, 충전전류 증가 램프 ≥ 600 µs(375 A 기준 ≤ 0.625 A/µs). 640–920 V·보조부하 15·25 kW에서 확인했고, 640 V·1 kW는 위상 첨두 78.74 A로 공칭 한도 74.53 A를 넘어(E46) 경부하를 포함한 충분조건이 아니다 & 도출(E16, TB07, E35, E46) \\
고장 강등 & 경로·필터 ≤ 100 A, 직렬·링크 ≤ 120 V, 바이패스 ≤ 329 A, 버스 ≥ 260 V, 400 V 부하 연결 중 버스 ≤ 상한, 강등 상태 도달 & E17·TB09 \\
강등 중 DAB 첨두 & 링크 쪽 순시 첨두 ≤ 1.5 × 정격점 첨두(R2-W460 68.7 A, R2-W410 90.0 A) & 도출(E20) \\
보호 하드웨어 & 링크 제동 저항 3 Ω(제동 에너지 최대 51.9 J의 펄스열. 펄스열의 열누적·반복 주기 확인을 조건으로 둔 후보), 바이패스 게이트 선형 영역 구동, 400 V 부하 분리 경로 & 도출(E17, E40) \\
Stage C 보호 일시 정지 & 링크 붕괴·소프트 투입 동안 4 kW/ms로 입력을 0으로, 최대 12.4 ms. 그동안 48 V 부하는 48 V 배터리가 공급하고 Stage C의 ±2 \% 조정은 적용되지 않는다(48 V = 배터리 전압). 에너지는 포트별: 고압 쪽 Stage C 입력 사각 상계 5.3 kW × 12.41 ms = 65.8 J(램프 결손 적분 58.8 J, E38), 48 V 쪽 별도 시험 +5 kW 부하 12.4 ms = 62 J·배터리 등가(20 mΩ 세 점) 방전 64.7–65.3 J·41.6–45.7 V(E43). 배터리 요구: 36 V 유지에 OCV 44 V에서 내부저항 ≤ 57.6 mΩ, 방전 ≥ 139 A(12.4 ms) & 도출(E38, E43) \\
이중 고장 400 V 부하 정지 & 전용 배선 신호 뒤 3 ms 안에 0(방전 저항 없음) & 도출(E38) \\
48 V 회생 차단 & 고장 검출 뒤 R2-W460 ≤ 1.0 ms·R2-W410 ≤ 0.5 ms 안에 시작, 4 kW/ms & 도출(E29) \\
400 V급 급속충전 / 48·12 V 전원 & 충전기 명령 I = min[375 A, 150 kW/(V\textsubscript{bat}/2)](640·800·920 V에서 120·150·150 kW, 920 V 약 326.1 A). 손실 비교(E38)는 375 A 고정(120·150·172.5 kW). 충전기 직결 단자는 V\textsubscript{A}(셀 3.7 V·375 A·보조부하 0의 단자 출력 125.19–166.82 kW, 180–240s)이고 팩별 손실 비교의 150 kW는 (V\textsubscript{bat}/2)I 기준 / 48 V 3 kW + 12 V 1.5 kW 동시(Stage C 5 kW; 12 V는 48 V 파생 부하, 독립 12 V 단자 조정 미시험), 48 V 회생 −3 kW & 부하 등급 \\
조정형 Stage C & 입력 270–504 V, 48 V 설정에서 ±2 \%(계약 과도 ≤ 4 kW/ms)·36–52 V(계단 견고성). 36–52 V CV 설정 범위 전체는 보장하지 않음(52 V·270 V·−650 W 단측 송신 듀티 0.980 > 0.97, 10.2절, E46), 출력 정격 105 A(과도 130 A, 2 ms), 5 kW 효율 ≥ 96 \% & 규격\cite{b10}, E22 \\
Stage C 데드타임 & 1차 100 ns·2차 40 ns(두 브리지 구동 때 90 ns). 출력용량 ±20 \%에서 통과하는 창의 가운데이며 창 폭은 20–30 ns & 도출(E22, TB14·TB15) \\
Stage A 영전류 스위칭 & 턴오프 전류 ≤ max(5 \% × 반주기 첨두, 1 A), 위상 첨두 ≤ 정격 충전점 첨두(공칭 74.5 A), 모듈 1개 고장 시 ≤ 1.5배(111.8 A) & E1, E24, E26, E35 \\
Stage A 위상 셰딩·모듈 관리기 & 위상: 멈춤 20 A(600 µs), 복귀 30 A 또는 순시 53 A, 자연 턴오프 뒤 멈춤. 모듈: 추가 70 A/모듈·순시 66 A·충전 요청, 줄임 60 A/모듈·1 ms, 고장 모듈 영구 제외 & 도출(E24, E28, E30, E35) \\
충전기 비상 정지 & (CHAdeMO 방식, IEEE 2030.1.1-2015 부속서 A 시험 항목) 충전 허가 신호가 꺼진 뒤 30 ms 안에 출력 ≤ 5 A(이 보고는 30 ms 동안 정격 전류 유지를 보호 포락선으로 씀), 차량 접촉기는 5 A 이하에서 개방, 접촉기 차단 시 커넥터 ≤ 600 V(서지 제외). 다른 충전 방식은 확인하지 않음 & CHAdeMO 시험 항목\cite{b16} \\
V\textsubscript{A} 클램프 & 468 V 켬·463 V 끔(V\textsubscript{bat} 최대 920 V 이하 팩), 224s·240s는 511/506 V. 가지 3개 × 891SP 2.4 Ω(W460 2.7 Ω) + 가지마다 SiC 스위치, 스위치 켬·끔 지연 목표 각각 ≤ 1 µs(검증 입력은 각각 1.0 µs 고정, 0.5 µs 열 판정 반례 E45, 12.6절), 최대 연속 켜짐 0.5 ms, V\textsubscript{A} 사건 한도 470·530 V. 문턱 구현 요구는 표 \ref{t:e26r} & 도출(E26, E38–E40, E44, E45, TB20, TB28–TB33) \\
Stage A 공진주파수·출력용량 & 영전류 게이트에서 스위치 출력용량의 하드 스위칭 손실(전이당 Q\textsubscript{oss}·V)을 넣어 판정. 조사한 후보 90개 중 150·100 kHz의 30개가 영전류 게이트 조건의 E31 동시 기준에서 0/30이어서 25 kHz, V\textsubscript{A} 노드 커패시터 2.4 mF. 공진 인덕터는 공심 & 도출(E31–E35) \\
부품 정격 여유 & 전압 ≤ 0.8 × 정격, 전류 ≤ 0.5 × 펄스 정격, 단발 T\textsubscript{J} ≤ T\textsubscript{J,max}, 정상 운전 T\textsubscript{J} ≤ T\textsubscript{J,max} − 25 °C. 열저항·손실 인자는 데이터시트 최대값, 전형값만 있으면 규칙 L(16장) & E34, E37–E40 \\
\end{longtable}
}
\srcnote{주: 근거 번호는 부록 B의 검증 근거를 따른다.}

\section{가설별 수락 기준}

가설은 요구사항에서 도출해 시험 전에 고정한 기준으로 판정한다(표 \ref{t:accept}). H5는 처음 기준 그대로 판정했고, 그 결과(기각)에 따라 적용 범위를 묻는 판정을 더했다. 결과는 18.1절에 정리한다.

{\footnotesize
\begin{longtable}{L{\dimexpr 0.2105\linewidth-2\tabcolsep\relax}L{\dimexpr 0.3509\linewidth-2\tabcolsep\relax}L{\dimexpr 0.4386\linewidth-2\tabcolsep\relax}}
\caption{가설별 판정 지표와 수락 기준}\label{t:accept}\\
\toprule
\tabhead \textbf{가설(질문)} & \textbf{판정 지표} & \textbf{수락 기준} \\
\midrule
\endfirsthead
\toprule
\tabhead \textbf{가설(질문)} & \textbf{판정 지표} & \textbf{수락 기준} \\
\midrule
\endhead
\bottomrule
\endlastfoot
H1 정격 분리(Q1) & Stage B 링크전력, 충전전류 경로 & 링크 ≤ 8 kW(정상·Tier-1·Tier-2, 공차 포함), 충전전류 Stage B 비경유 \\
H2 강건 제어(Q2) & Tier-1·Tier-2, 파라미터 ±10 \% 공차, 지연 예산 & Tier-1: 전압 품질 + 전류 기준. Tier-2: 물리정격 + 전류 기준 + 버스 창. 공차와 추가 지연 3.0 µs에서 Tier-1 유지 \\
H3 충전 중 보조버스(Q1) & 150 kW(375 A 제한) 충전 + 25 kW 보조부하, 충전기 접속·해제(충전기 명령 I = min[375 A, 150 kW/(V\textsubscript{bat}/2)]은 입력 한계이며 세 전압에서 배터리 순유입 150 kW 유지를 뜻하지 않음) & 동시 운전에서 Tier-1, 모드 전환에서 Tier-2 \\
H4 부하 승인(Q2) & 부하 계약, 메시지 지연 1–5 ms, 손실 0–20 \% & 예산 초과 0, 컴프레서 승인 지연: 메시지 무손실 최대 ≤ 25 ms, 손실 ≤ 20 \%에서 99백분위 ≤ 50 ms, 미공급이 정적 예약보다 작을 것 \\
H5 기준 구조 대비 이점(Q3) & 기준 구조 ①·②·④와 변형 ③′ 대비 동일 경계 비교 & 손실·정격 VA·반도체 수 중 2개 이상에서 ④와 ③′보다 우수, 서비스 동등 이상 \\
Q3 적용 범위(추가) & 셀 180–240s × 부품 B\textsubscript{420}·B\textsubscript{470}의 성립성과 손실 & 성립 구조 가운데 최소 손실인 팩·부품 조합이 있을 것 \\
Q3 산업 방식 대비(추가) & 모터·인버터 승압 충전(⑤), 800 V급 보조부품 설계(⑥) 대비 손실·반도체 수 & R2가 앞서는 조건과 뒤지는 조건을 같은 손실 규칙으로 구분할 것 \\
\end{longtable}
}
\srcnote{주: 구조 번호 ①–⑥은 부록 A의 약어 표를 따른다.}
